\chapter{相关理论与技术}

\section{深度自编码器原理}

自编码器是一种无监督神经网络，由编码器和解码器两部分组成。编码器将输入向量映射为低维隐表示，解码器再由隐表示重构输入。设输入向量为 $x$，编码器输出为 $h$，则编码与解码过程可表示为
\begin{equation}
  h = f(W_1 x + b_1),
  \label{eq:ch2-encoder}
\end{equation}
\begin{equation}
  \hat{x} = g(W_2 h + b_2),
  \label{eq:ch2-decoder}
\end{equation}
其中 $W_1$、$W_2$ 为权重矩阵，$b_1$、$b_2$ 为偏置向量，$f$ 与 $g$ 为激活函数。模型以输入 $x$ 与重构结果 $\hat{x}$ 之间的差异作为优化目标，如式 \eqref{eq:ch2-decoder} 所示。训练完成后，重构误差较大的样本被判定为异常。

\section{注意力机制}

注意力机制能够自适应地为不同特征分配权重，使模型重点关注对检测结果贡献更大的特征维度。本文将注意力层嵌入自编码器的编码阶段，提升了模型对关键流量特征的敏感性。

\section{评价指标}

本文采用准确率、召回率和 $F_1$ 值作为评价指标。设真正例、假正例、真负例、假负例分别为 $TP$、$FP$、$TN$、$FN$，则
\begin{equation}
  F_1 = \frac{2TP}{2TP + FP + FN}.
  \label{eq:ch2-f1}
\end{equation}

\subsection{准确率与召回率}

准确率反映检测结果中正确样本所占比例，召回率反映真实异常样本中被正确检出的比例，两者从不同侧面描述模型性能。

\subsection{$F_1$ 值}

$F_1$ 值是准确率与召回率的调和平均，能够综合衡量模型性能，适用于正负样本不均衡的检测场景。
